% Generated draft: requires mathematical and source review.
\documentclass[11pt]{article}
\usepackage[margin=25mm]{geometry}
\usepackage{fontspec}
\setmainfont{Latin Modern Roman}
\usepackage{amsmath,amssymb,mathtools,unicode-math}
\setmathfont{Latin Modern Math}
\usepackage{enumitem,needspace,xurl,hyperref}
\setlength{\parindent}{0pt}
\setlength{\parskip}{5pt}
\setlength{\emergencystretch}{3em}
\newcommand{\N}{\mathbb{N}}
\newcommand{\Z}{\mathbb{Z}}
\newcommand{\Q}{\mathbb{Q}}
\newcommand{\R}{\mathbb{R}}
\newcommand{\C}{\mathbb{C}}
\newcommand{\F}{\mathbb{F}}
\newcommand{\A}{\mathbb{A}}
\newcommand{\PP}{\mathbb{P}}
\newcommand{\E}{\mathbb{E}}
\newcommand{\eps}{\varepsilon}
\newcommand{\dd}{\,\mathrm{d}}
\begin{document}
\section*{UnsolvedMath ID: 30004566}
\textbf{Source problem code:} catalogue entry
\subsection{Definitions and mathematical statement}
The source asks how knot ropelength depends on crossing number, without specifying the ambient space, regularity, thickness convention, or a particular conjectured inequality. The following definitions make explicit the standard classical-knot interpretation; these conventions require confirmation against the source report.

Let \(K\) be a tame knot type in \(\R^3\), considered up to ambient isotopy. For a regular embedded \(C^{1,1}\) closed curve \(\gamma:S^1\to\R^3\) representing \(K\), let \(L(\gamma)\) denote its length and define its thickness by
\[
\operatorname{Thi}(\gamma)=\sup\bigl\{r>0: \text{every }x\in\R^3\text{ with }\operatorname{dist}(x,\gamma(S^1))<r\text{ has a unique nearest point on }\gamma(S^1)\bigr\}.
\]
Define the ropelength of the knot type by
\[
\operatorname{Rop}(K)=\inf_{\substack{\gamma\text{ represents }K\\ \gamma\text{ regular, embedded, }C^{1,1}\\ \operatorname{Thi}(\gamma)>0}}
\frac{L(\gamma)}{\operatorname{Thi}(\gamma)}.
\]
Let \(\operatorname{Cr}(K)\) be the minimum number of crossings among all regular planar diagrams representing \(K\), where crossings are transverse double points supplied with overcrossing and undercrossing information.

\textbf{Source question, expressed as a quantitative research problem.} Determine sharp uniform relationships between \(\operatorname{Rop}(K)\) and \(\operatorname{Cr}(K)\). In particular, seek functions \(f,g:\{0,1,2,\ldots\}\to[0,\infty)\) for which
\[
 f\bigl(\operatorname{Cr}(K)\bigr)\leq \operatorname{Rop}(K)\leq g\bigl(\operatorname{Cr}(K)\bigr)
\]
for every knot type \(K\) in the intended class, and determine the optimal possible asymptotic orders of such bounds as crossing number tends to infinity. Under the classical interpretation above, the class consists of all tame knot types and includes the unknot, whose crossing number is zero. Bounds restricted to nontrivial knots or to sufficiently large crossing number would require an explicit restriction.

The request for sharp uniform bounds is an interpretation of the broad source question, not a supplied specific conjecture. The record does not prescribe particular functions, exponents, logarithmic factors, a meaning of sharpness, or a restriction to a subclass of knots. These intended targets and conventions remain unresolved.
\subsection{Short English statement}
Ropelength measures the least possible length of a knot relative to its thickness; crossing number measures the fewest crossings in any diagram of its knot type. The question asks for sharp quantitative bounds relating these invariants. The supplied formulation does not specify a particular inequality or asymptotic conjecture, and the classical-knot and thickness conventions used here require source confirmation.
\subsection{Sources}
\begin{itemize}
\item Local UnsolvedMath registry, supplied complete source record: permanent ID \(30004566\), source problem code \texttt{catalogue entry}, title \emph{Sharp Relations Between Knot Ropelength and Crossing Number}, collection \emph{Oberwolfach Reports Open Problems}. No individual proposer is identified; the supplied proposed year is 2020.
\item Supplied primary report citation: \emph{Mini-Workshop: Nonlocal Analysis and the Geometry of Embeddings} (2021), \emph{Oberwolfach Reports}, pp. 1823--1856, DOI \texttt{10.4171/owr/2020/37}. \url{https://doi.org/10.4171/owr/2020/37}. The record attributes the problem to workshop open-problem discussions but gives no problem-specific page.
\item Supplied literature reference: Diao et al., ropelength and crossing number. \url{https://arxiv.org/abs/0912.3282}.
\item Supplied literature reference: Cantarella, Faber and Mullikin, upper bounds for ropelength. \url{https://doi.org/10.1016/S0166-8641(03)00168-8}.
\item Supplied literature reference: Buck and Simon, thickness and crossing number. \url{https://doi.org/10.1016/S0166-8641(97)00211-3}.
\item Supplied dated literature curation, checked \texttt{2026-08-22}: status \texttt{partially\_solved}. Its assessment reports substantial upper and lower bounds, including near-linear upper bounds with logarithmic loss, while describing the sharp asymptotic relation and optimal exponents as unknown. This is supplied dated curation derived from the OWR record, not an independently verified literature assessment; the status and cited claims require independent verification.
\item The supplied statement status is \texttt{corrected\_verified}. Its accompanying verification note says that only formula delimiters, notation, grammar, and typography were repaired, preserving mathematical content and scope. This label does not establish independent mathematical verification of this draft.
\end{itemize}
\end{document}
